\section{The conditional record and the Schinzel audit}
\label{sec:record}

\subsection{The statement, with its hypothesis exposed}

The member property below records the output required by L6; it is no
longer an additional conjectural premise.

\begin{definition}[Member property $\mathcal M$]
For every odd prime $w$ and every rational $z$ with $v_w(z)\geq1$, there
is a verified aligned class for $(w,z)$ and a prime member
$Q\equiv q\pmod N$ outside its finite excluded set whose bad emergent set
is empty.  Explicitly, for every odd prime
$\ell\notin S\cup\{Q\}$ with
$v_\ell(P(\varepsilon fQ))$ odd, the local Hilbert symbol
$(x_0,d_0)_\ell$ is $+1$, where $S$ is the frozen set of the class.
\end{definition}
% src: data/l19_classexist.jsonl (uniform class component)
% src: data/l18_schinzel_implies_h.jsonl (member contract)

\begin{theorem}[CONDITIONAL ON CLASSICAL SCHINZEL H: Theorem C]
Assume classical Schinzel's Hypothesis~H.  For every odd cell, use the
fixed canonical branch
\[
 \tau^\dagger=\frac{1+2a^2}{1+4a^2},
 \qquad \delta_{\tau^\dagger}=-\frac{4a^4}{1+4a^2}.
\]
Then the member property $\mathcal M$ holds and
$\Q\setminus\Z$ is Diophantine in $\Q$ with six unknowns,
realized by the L6 formula
\[
 F(z):=\exists b,s\;\bigl[(1+2s,b)\in\Phi_1^{\{2\}}\ \wedge\
 \Theta_*\bigl(1+2s,b,h(1+2s,b,z^3),s\bigr)\bigr],
\]
where, with $A=1+4a^2$, $B=2b$, and
$\tau^\dagger=(1+2a^2)/A$,
\[
\begin{aligned}
 \Theta_*(a,b,c,s):=\exists y,r\quad&
 [c^2-Ay^2-16Br^2=16-16ABs^2]\\
 &{}\vee[c^2-Ay^2-16ABr^2=16A-16A^2Bs^2]\\
 &{}\vee[a^4(c^2-Ay^2)+4ABr^2=4A^2Bs^2-4A].
\end{aligned}
\]
The three disjuncts remain the cleared equations for
$\tau=0$, $\tau=2a/A$, and $\tau^\dagger$, respectively, in the same
witnesses $(y,r)$.  They preserve the finite soundness menu and the
six-count; completeness in this theorem uses only the third, fixed
canonical disjunct.
% src: data/l19_classexist.jsonl (square-branch definition)
% src: data/l22_elimination.jsonl (fixed-Z-theorem row)
% src: CONDITIONAL.md:18-31
Consequently $\Z$ is definable in $\Q$ by a formula with six universal
quantifiers,
$\Z=\Q\setminus(F\vee\mathfrak m_2)$, and
$\operatorname{efd}_{\Q}(\Q\setminus\Z)\leq 5$.
% src: CONDITIONAL.md:14-25
\end{theorem}

This is a conditional definition, not an unconditional solution of
Hilbert's Tenth Problem over $\Q$.  Soundness of the finite branch menu
consumes no hypothesis.  For completeness, fixed-$Z$ vertical
irreducibility and Hilbert selection give a good canonical-branch $a$ in
the existing admissible progression; L20 supplies the class and all
Schinzel-pair conditions; Theorem~\ref{thm:schinzel} turns classical
Schinzel~H into an emergent-free member; and L6 gives the displayed
definition.  No per-cell irreducibility certificate is assumed.
% src: data/l22_elimination.jsonl (fixed-Z-theorem row)
% src: data/l20_admissible.jsonl; data/l18_schinzel_implies_h.jsonl

The quantifier saving is attributed precisely to Nicolas Daans, Philip
Dittmann, and Arno Fehm, \emph{Existential rank and essential dimension of
diophantine sets}, arXiv:2102.06941v5, Theorem~1.4 (equivalently
Corollary~5.11)~\cite{ddf}.  Their intersection theorem turns the ranks
three and two on the common base into $3+2-1=4$ witnesses; the two base
coordinates then give six unknowns.  The recorded rank--essential-dimension
identity gives the bound five.  These counts remain \textbf{CONDITIONAL}
because the set equality to $\Q\setminus\Z$ uses classical Schinzel~H.
% src: CONDITIONAL.md:31-48
% src: data/litscout_h10q.md:30-34

\subsection{What is proved, and what Schinzel must supply}

\textbf{PROVED (uniform class side and linear side).}
Theorem~\ref{thm:classexist} constructs an aligned class for every cell and
freezes all controlled Hilbert symbols while restricting the moving prime to
one reduced progression
\[
 q_1(t)=q_{1,0}+Nt,\qquad \gcd(q_{1,0},N)=1.
\]
The residue construction gives infinitely many proven-prime choices of the
base $q_{1,0}$, and Dirichlet's theorem supplies infinitely many prime values
of the progression.  This is the linear, unconditional half of the
member-existence problem; no probable-prime test or refusal is used as
evidence.
% src: data/l19_classexist.jsonl (meta, construction rows, and summary)

\textbf{CONDITIONAL (the sole conjectural cofactor side).}  Put
\[
 F_{\mathrm{cell}}(t)=P\bigl(\varepsilon f q_1(t)\bigr),
\]
using the primitive integer representative after removal of the fixed
rational scale recorded in the audit.  The L22 theorem and quantitative
Hilbert irreducibility make $F_{\mathrm{cell}}$ irreducible for a selected
admissible $a$.  Requiring it to be prime simultaneously with
$q_1(t)$ is the two-polynomial instance of classical Schinzel's
Hypothesis~H \cite{schinzel}.  The polynomial has degree eight.  A prime
cofactor gives an emergent-free member because the class has frozen the
controlled places and reciprocity forces the remaining single symbol.
After using the progression prime as variable, the pair is
$\{q_1+Nt,F_{\mathrm{cell}}(t)\}$.
% src: data/l22_elimination.jsonl (fixed-Z-theorem row)
% src: data/l20_admissible.jsonl; data/l18_schinzel_implies_h.jsonl
% src: THEOREMS.md:953-971
% src: CONDITIONAL.md:175-185
% src: data/l17_schinzel_audit.jsonl:1-294

\subsection{PROVED finite audit of the Schinzel pair}

The exact audit contains $293$ canonical records, split into $190$ L11
classes and $103$ escape classes.  The fields
\texttt{F\_irreducible}, \texttt{F\_value\_gcd}, and
\texttt{pair\_product\_gcd} respectively certify irreducibility on
$293/293$ rows, value gcd one on $293/293$ rows, and pair-product gcd one
on $293/293$ rows.  Every irreducibility certificate is a degree-eight
Frobenius certificate modulo a certified prime; both gcds were computed
exactly on $0\leq t\leq 2000$.  The pair-product computation is already a
certificate that no prime divides the product for every integer $t$: any
such fixed divisor would divide the computed gcd.  Thus all canonical
pairs pass the fixed-divisor and irreducibility conditions.  This is
\textbf{PROVED} for the stored pairs, not a prime-value theorem.

The finite audit does not by itself extrapolate to arbitrary rational
cells.  The L20 audit proves normalization, coprimality, absence of a fixed
divisor, frozen-unit and symbol constancy, and positive leading coefficient
uniformly.  Theorem~\ref{thm:fiber-irred} now supplies the missing uniform
input directly: for every fixed nonzero rational $Z$, $P(a,Z)$ is
irreducible in $\Q(a)[b]$.  Corollary~\ref{cor:admissible-irred} then
selects a good $a$ inside the character progression.  Thus the
$353$ recorded target certificates remain exact finite corroboration, but
they are no longer a premise for a new cell.
% src: data/l20_admissible.jsonl; data/l21_irred_generic.jsonl
% src: data/l22_elimination.jsonl (fixed-Z-theorem row)
% src: data/l17_schinzel_audit.jsonl:1-294

For a prime $p$ dividing $N$, the corrected residue count is
\[
 \nu_p=\#\{t\in\F_p:q_1(t)\not\equiv0\pmod p,
                  \ F_{\mathrm{cell}}(t)\not\equiv0\pmod p\}.
\]
The first condition is essential.  Since $p\mid N$ and
$\gcd(q_{1,0},N)=1$, the linear polynomial is the nonzero constant
$q_{1,0}$ modulo $p$; it must not be replaced by the parameter $t$.
Consequently the count is $p$, rather than $p-1$, whenever the cofactor has
no root modulo $p$.  The audit enumerated every $p\leq1000$ dividing the
row's modulus and obtained exactly this full count in every occurrence.
% src: data/l17_schinzel_audit.jsonl:1-294

\begin{table}[ht]
\centering
\small
\begin{tabular}{c|c|c|c}
primes $p$ & records containing $p\mid N$ & range of $\nu_p$ & mean density $\nu_p/p$\\ \hline
$2,3,5,7$ & $293$ each & $p$ & $1$\\
$11,37,89,97$ & $13$ each & $p$ & $1$\\
$13,53$ & $26$ each & $p$ & $1$\\
$17$ & $39$ & $p$ & $1$\\
$23$ & $38$ & $p$ & $1$\\
$29,41,61,73$ & $28$ each & $p$ & $1$\\
$43,47,67,83$ & $10$ each & $p$ & $1$
\end{tabular}
\caption{Corrected admissible-residue counts for primes dividing the progression modulus.}
% src: data/l17_schinzel_audit.jsonl:1
\end{table}

These local counts neither assume independence nor estimate a density of
simultaneous prime values.  Their role is narrower and exact: together with
the pair-product gcd, they rule out the elementary local obstruction for
each stored Schinzel pair.
% src: data/l17_schinzel_audit.jsonl:1

\subsection{Residual certificates and the literature boundary}

The six residual rows are deliberately segregated.  They are rational or
composite-$b$ witnesses verified by exact branch reconstruction and tied
status; they are not coerced into the prime family and are not evidence for
Schinzel's hypothesis.

\begin{table}[ht]
\centering
\small
\begin{tabular}{c|c|c|c}
cell $(w,z)$ & $(a,b)$ & $\tau$ & status\\ \hline
$(67,2)$ & $(1,4757)$ & $3/5$ & PROVED pointwise\\
% src: data/l17_schinzel_audit.jsonl:295
$(41,1/7)$ & $(25,4985)$ & $1251/2501$ & PROVED pointwise\\
% src: data/l17_schinzel_audit.jsonl:296
$(61,-2/7)$ & $(25,55571)$ & $1251/2501$ & PROVED pointwise\\
% src: data/l17_schinzel_audit.jsonl:297
$(61,2)$ & $(25,-71431)$ & $1251/2501$ & PROVED pointwise\\
% src: data/l17_schinzel_audit.jsonl:298
$(53,1/3)$ & $(-15,9)$ & $451/901$ & PROVED pointwise\\
% src: data/l17_schinzel_audit.jsonl:299
$(29,-2)$ & $(-15,-29/5)$ & $4055/2703$ & PROVED pointwise
% src: data/l17_schinzel_audit.jsonl:300
\end{tabular}
\caption{The six residual witnesses: unconditional certificates for their individual cells only.}
\end{table}

\textbf{OPEN (the unconditional published boundary).}  The exact
square-class reference is David Krumm, \emph{Squarefree parts of polynomial
values}, Journal de Th\'eorie des Nombres de Bordeaux \textbf{28} (2016),
no.~3, 699--724, arXiv:1407.4890~\cite{krumm}.  Its Theorem~1.3 and
Propositions~3.4--3.5 give the unconditional degree-one and degree-two
cases; Proposition~3.8 makes the degree-three case conditional on the
Parity Conjecture for elliptic curves over $\Q$.
% src: data/litscout_h10q.md:6-13

At prime arguments, the nearest unconditional results reach the cubic
boundary without prescribing a square class: Helfgott proves square-free
values of $f(p)$ for separable cubic $f$ \cite{helfgott}, while Reuss
proves $(d-1)$-free values rather than prime or prescribed-square-class
values \cite{reuss}.  Neither yields the required degree-eight
conclusion.
% src: data/litscout_h10q.md:15-28

The conic-bundle literature does not remove the gap.  The unconditional
Harpaz--Skorobogatov--Wittenberg theorem requires abelian constant fields
in the bad fibres \cite[Theorem~3.1 and Corollary~3.4]{hsw}.  The general
Schinzel reference in \emph{J. reine angew. Math.} \textbf{495} (1998),
1--28 is by Colliot-Th\'el\`ene, Skorobogatov and Swinnerton-Dyer
\cite{ctssd}, not by Colliot-Th\'el\`ene and Sansuc; the latter pair's
Schinzel paper is \cite{ctsansuc}.  No cited unconditional theorem supplies
simultaneous prime values for the present linear--octic pair.  Classical
Schinzel~H is therefore the sole conjectural input to the six-quantifier
theorem, which remains \textbf{CONDITIONAL}; no unconditional record or
solution of H10($\Q$) is claimed.
% src: data/litscout_h10q.md:36-39
